\newpage\SourceLeaf{44}\ArithmeticHead{37}{R}{KREISTEILUNG.}
\begin{center}\bfseries 15.\end{center}
\noindent\hspace*{\parindent}\textsc{Gauss} hat aber ferner jede der Hilfsgleichungen auf eine reine Gleichung zurückzuführen gelehrt, indem er sich dazu der Resolvente von \textsc{Lagrange} bediente. Die Verwendung der Resolvente ist in \textsc{Gauss}’ \textit{Tagebuch} (Nr. 65, 66 T.) auf den 17. Juli 1797 datiert, denn die dort angegebene deductio secunda ist nichts anderes, als die in art. 360 der D. A. gelehrte Methode, wie aus dem Passus: »quae theoriam secundam aequationum purarum in art. 360 D. A. inchoatam magis illustrant« (W. II, S. 263, Ende von Nr. 18) deutlich hervorgeht, während mit der Notiz Nr. 55 T. vom 19. Januar 1797 die \textls[100]{erste} Methode der sukzessiven Hilfsgleichungen gemeint sein dürfte.

Übrigens tritt der Gedanke an das Mittel der Resolvente schon früher (am 17. Sept. 1796, Nr. 37 T.) anscheinend selbständig bei \textsc{Gauss} auf, wie er denn auch andere, elementarere Stücke der Lehre von den Gleichungen, wie die \textsc{Newton}schen Formeln (Nr. 6 T. 23. Mai 1796; Nr. 28 T. 21. Aug. 1796) die nebst ihrer Umkehrung sich in \textsc{Leiste}, \textit{Arithm. u. Algebra} zu S. 6 u. 7 notiert finden (siehe W. X 1, S. 127), und Anderes für die Entwicklung seiner Kreisteilungstheorie sich erst selbst zurechtgelegt zu haben scheint. Setzt man nun, unter \(R\) eine primitive \(e^{\mathrm{te}}\) Einheitswurzel verstehend,
\[
\begin{gathered}
[r,R^i]=r+R^i r^g+R^{2i}r^{g^2}+\ldots+R^{(p-2)i}r^{g^{p-2}},\\
(i=1,2,\ldots,e-1)
\end{gathered}
\]
wofür auch
\[
\begin{gathered}
[r,R^i]=\eta_0+R^i\eta_1+R^{2i}\eta_2+\ldots+R^{(e-1)i}\eta_{e-1}\\
(i=1,2,\ldots,e-1)
\end{gathered}
\]
geschrieben werden kann, so findet sich
\[
[r,R^i]^e=T_i,
\]
\newpage\SourceLeaf{45}\ArithmeticHead{38}{L}{BACHMANN, ÜBER GAUSS’ ZAHLENTHEORETISCHE ARBEITEN.}
\noindent wo \(T_i\) gleichzeitig mit der Wurzel \(R\) als bekannt betrachtet werden kann. Durch Auflösung dieser \(e-1\) reinen Gleichungen \(e^{\mathrm{ten}}\) Grades und mit Zuhilfenahme der Gleichung \([r,1]=-1\) findet sich dann sogleich
\[
\eta_0=\frac1e\bigl(-1+\sqrt[e]{T_1}+\sqrt[e]{T_2}+\ldots+\sqrt[e]{T_{e-1}}\bigr),
\]
d. i. die Auflösung der Gleichung \(X_e=0\), und ähnlicherweise werden die übrigen Hilfsgleichungen gelöst. Die allzugrosse Vieldeutigkeit dieser Formel beschränkt man auf das zutreffende Mass, wenn man beachtet, dass allgemein auch der Ausdruck
\[
[r,R^i]\mathbin{.}[r,R]^{e-i}=T^{(i)}
\]
eine bekannte Grösse ist, sodass die vorige Formel auch in die Gestalt
\[
\eta_0=\frac1e\biggl(-1+\sqrt[e]{T_1}+\frac{T^{(2)}}{T_1}\mathbin{.}(\sqrt[e]{T_1})^2+\ldots+\frac{T^{(e-1)}}{T_1}\mathbin{.}(\sqrt[e]{T_1})^{e-1}\biggr)
\]
gesetzt werden kann, aus welcher dann die übrigen Perioden \(\eta_1,\eta_2,\ldots,\eta_{e-1}\) hervorgehen, wenn der einen noch auftretenden Wurzel ihre sämtlichen Werte beigelegt werden. Nur irrtümlicherweise konnte, wie \textsc{Gauss} bemerkt hat (W. II, S. 249), \textsc{Lagrange}\footnote{\textit{Traité de la résolution des équations numériques}, 2. éd, 1808, Note XIV, art. 41, Oeuvres VIII, S. 367.}\ die erstere Formel als die vorzüglichere bezeichnen. Der hierbei auftretende Zweifel, ob \(T_1=[r,R]\) auch nicht Null sei, kann, wie schon \textsc{Gauss} bemerkt, aber nicht weiter bewiesen hat, behoben werden. Auch gab schon \textsc{Gauss} (W. II, S. 252) in einer Abhandlung, welche aus dem Jahre 1808 stammt (W. II, S. 265, Bemerkungen), den später von \textsc{Jacobi} u. A. hergeleiteten Summen-Ausdruck für den Quotienten
\[
\frac{[r,R^i]\mathbin{.}[r,R^k]}{[r,R^{i+k}]},
\]
insbesondere auch die Formel
\[
[r,R^i]\mathbin{.}[r,R^{-i}]=(-1)^i\mathbin{.}p
\]
und zeigte (W. II, S. 250), dass zur Bestimmung von \(\sqrt[e]{T_1}\) bei Hinzunahme der \(e^{\mathrm{ten}}\) Einheitswurzeln, d. h. neben der Teilung der ganzen Peripherie in \(e\) gleiche Teile, die Teilung eines gegebenen Winkels in ebensoviel gleiche Teile
\newpage\SourceLeaf{46}\ArithmeticHead{39}{R}{KREISTEILUNG.}
\noindent nebst der Ausziehung der Quadratwurzel aus einer bekannten Grösse genügt. Insonderheit erfordert die Auflösung der Kreisteilungsgleichung \(X=0\) selbst, d. h. die Teilung des Kreises in \(p\) gleiche Teile, nur die Teilung des Kreises und die eines dann gegebenen Winkels in \(p-1\) gleiche Teile nebst der Ausziehung der Quadratwurzel aus einer bekannten Grösse, nämlich aus \(p\).

Die Wurzeln von \(X=0\), welche einer der \(f\)-gliedrigen Perioden \(\eta_i\) angehören, sind Wurzeln einer andern Gleichung vom Grade \(f\), deren Koeffizienten linear durch die Perioden \(\eta_i\) ausdrückbar sind. Daher zerfällt, wenn die letzteren bekannt geworden sind, der Ausdruck \(X\) in \(e\) Faktoren \(f^{\mathrm{ten}}\) Grades mit rational bekannten Koeffizienten. Als Gleichung für die zwei \(\frac{p-1}{2}\)-gliedrigen Perioden fand \textsc{Gauss} die folgende:
\[
x^2+x+\tfrac12\bigl(1-(-1)^{\frac{p-1}{2}}p\bigr)=0
\]
und die zugehörige Zerlegung
\[
4X=Y^2-(-1)^{\frac{p-1}{2}}pZ^2,
\]
wo \(Y,Z\) ganze ganzzahlige Funktionen von \(x\) sind. Auf jener Gleichung beruhen die beiden Beweise des Reziprozitätsgesetzes, welche die A. R. (W. II, S. 234) enthält; den \textsc{Gauss}schen Tagebuchnotizen zufolge (s. Nr. 30 T.) muss sie ihm also schon vor dem 2. September 1796 bekannt gewesen sein. Der Herleitung dieser Gleichung entnahm er aber auch (D. A. art. 356) einen neuen Beweis für den quadratischen Charakter \(\left(\frac{-1}{p}\right)\) und ward hier schon auf die sogenannten \textsc{Gauss}schen Summen geführt, die uns noch nachher beschäftigen werden. Desgleichen fand er am 1. Oktober 1796 zunächst durch Induktion (Nr. 39 T.) die Gleichung für die drei \(\frac{p-1}{3}\)-gliedrigen Perioden\footnote{Diese Gleichung, auf deren Aufsuchung schon anderweitige Notizen, wie die in \textsc{Leiste}, \textit{Arithm. u. Algebra} zu S. 8 und in \textsc{Hellwig}, \textit{Anfangsgründe der Mathematik} (siehe W. X 1, S. 111) deuten, wird in einer in \textsc{Lamberts} \textit{Tabellen} bei S. 223 befindlichen Aufzeichnung von \textsc{Gauss} unter der Form
\[
z^3-3pz-p(9a-p-1)=0
\]
gegeben, wobei
\[
4p=3N^2+(9a-p-1)^2.
\]
}\ und versah sie (Nr. 67 T.) dann am 20. Juli 1797 mit dem Beweise, welchen art. 358 der D. A. enthält und in dessen Verlaufe nebenbei die Zerlegung
\newpage\SourceLeaf{47}\ArithmeticHead{40}{L}{BACHMANN, ÜBER GAUSS’ ZAHLENTHEORETISCHE ARBEITEN.}
\noindent \(4p=x^2+27y^2\) für die Primzahlen von der Form \(p=6n+1\) gefunden wird; nach Nr. 135 des \textit{Tagebuchs} hat er diesen Beweis am 10. Mai 1808 auf wesentlich einfachere Grundsätze zurückgeführt, ohne jedoch letztere näher zu bezeichnen. Ebensowenig findet sich in seinen Schriften die Zerlegung von \(X\) in vier den Wurzeln der \(\frac{p-1}{4}\)-gliedrigen Perioden entsprechende Faktoren, die er (nach Nr. 128 T.) im Jahre 1806 noch ausgeführt hat. Nur die Hilfsmittel zur Bildung der dazu erforderlichen Hilfsgleichung vierten Grades liegen in den Betrachtungen der artt. 15—20 der ersten Abhandlung über die biquadratischen Reste bereit, mit denen sie, wie \textsc{Gauss} selbst in art. 22 dort angemerkt hat, aufs engste verbunden ist.

Aus den in dieser Skizze angegebenen Daten ist ersichtlich, dass zur Zeit, als \textsc{Gauss} seine Entdeckung bezüglich des Siebzehnecks bekannt gab, also im Juni 1796, ihm, wie er selbst dabei ausgesagt hat, zu einer vollständigen, logisch festgefügten Theorie der Kreisteilung doch noch eine ganze Reihe wesentlicher Sätze fehlte. Kaum wird man daher fehl gehen, wenn man jene Entdeckung als einen glücklichen Wurf, weniger als Ergebnis denn als Quelle der allgemeinen Methode betrachtet. In der Tat schreibt \textsc{Gauss} am 6. Januar 1819 an \textsc{Gerling} (siehe W. X 1, S. 125), dass ihm schon in seinem ersten Semester alles, was sich auf die Verteilung der Wurzeln in zwei Perioden bezieht, bekannt gewesen sei, aber erst am frühen Morgen des 29. März 1796 sei es ihm geglückt, den Zusammenhang der Wurzeln unter einander, d. i. das arithmetische Prinzip ihrer zyklischen Anordnung zu erkennen; und nun führte ihn bei der Siebzehnteilung die wiederholte Zweiteilung der Perioden nicht nur zum glücklichen Ziele, sondern auch zur Erkenntnis der allgemeinen Methode, deren Erfolg er sicher genug übersah, um ungeachtet der noch vorhandenen Lücken die oben erwähnte Anzeige wagen zu können.
\begin{center}\rule{.18\textwidth}{.3pt}\end{center}
